% !TeX program = xelatex
% 使用方式：复制本文件和 sz-campus-beamer.sty、assets 文件夹，再修改课程信息和各页内容。
\documentclass[UTF8,aspectratio=169,11pt]{ctexbeamer}
\usepackage{sz-campus-beamer}

%==================== 课程信息 ====================
\renewcommand{\CourseTitle}{这里填写课题}
\renewcommand{\CourseSubtitle}{这里填写章节或课型}
\renewcommand{\TeacherName}{华光辉}
\renewcommand{\LessonDate}{2026.09.06}
\renewcommand{\EndText}{敬请批评指正}

\title{\CourseTitle}
\author{\TeacherName}
\date{\LessonDate}

\begin{document}

%==================== 封面 ====================
\CampusTitlePage

%==================== 章节过渡页 ====================
\CampusSectionPage[从情境进入核心问题]{第一部分\quad 情境与问题}

%==================== 普通内容页 ====================
\begin{frame}{任务一：核心问题}
  \begin{questionbox}[观察与思考]
    这里填写本页最重要的问题。先让学生独立思考，再交流理由。
  \end{questionbox}

  \vspace{0.2cm}
  \begin{itemize}
    \item 已知条件：这里填写条件。
    \item 研究对象：这里明确本节课研究什么。
    \item 需要回答：这里写出可检查的学习任务。
  \end{itemize}
\end{frame}

%==================== 概念建立页 ====================
\begin{frame}{概念建立}{从实例到数学表达}
  \begin{definitionbox}[定义]
    这里填写教材定义，并用\keypoint{关键词}标出对象、条件和结论。
  \end{definitionbox}

  \pause
  \begin{columns}[T,totalwidth=\textwidth]
    \column{0.48\textwidth}
      \textbf{正例}\par
      这里填写满足定义的例子。
    \column{0.48\textwidth}
      \textbf{反例}\par
      这里填写容易混淆但不满足定义的例子。
  \end{columns}
\end{frame}

%==================== 双栏讲解页 ====================
\begin{frame}{表示与解释}
  \begin{columns}[T,totalwidth=\textwidth]
    \column{0.52\textwidth}
      \begin{questionbox}[符号语言]
        \[
          f(x)=x^2-2x+1=(x-1)^2
        \]
      \end{questionbox}
      这里解释符号中的对象、条件和数量关系。

    \column{0.44\textwidth}
      \begin{definitionbox}[图形或表格区域]
        插入图片时可使用：

        \medskip
        {\footnotesize\texttt{\textbackslash includegraphics[width=}\\
        \texttt{\textbackslash linewidth]\{图片\}}}
      \end{definitionbox}
  \end{columns}
\end{frame}

%==================== 例题页 ====================
\begin{frame}{例题训练}
  \begin{examplebox}[例1]
    这里填写完整题目。题目出现后先给学生独立思考时间。
  \end{examplebox}

  \pause
  \answer{这里填写答案。}
  \method{写出关键依据，说明为什么选择这一方法。}

  \begin{notebox}[易错点]
    这里记录学生可能漏掉的条件、混淆的概念或推理断点。
  \end{notebox}
\end{frame}

%==================== 小结页 ====================
\begin{frame}{课堂小结}
  \begin{enumerate}
    \item 本节课研究的核心对象是什么？
    \item 定义中不能遗漏的条件是什么？
    \item 哪一道题能够检验是否真正理解？
  \end{enumerate}

  \vspace{0.25cm}
  \begin{questionbox}[离堂检测]
    这里放一道可以独立完成、能够直接检验目标是否达成的小题。
  \end{questionbox}
\end{frame}

%==================== 结束页 ====================
\CampusEndPage

\end{document}
